\documentclass[11pt,a4paper]{article}
\usepackage[utf8]{inputenc}
\usepackage[margin=1in]{geometry}
\usepackage{xcolor}
\usepackage{hyperref}
\usepackage{enumitem}
\usepackage{fancyhdr}
\usepackage{titlesec}
\usepackage{booktabs}
\usepackage{microtype}
\usepackage{amsmath}

% Define brand color scheme
\definecolor{primary}{RGB}{20, 48, 82}     % Deep Royal Navy
\definecolor{secondary}{RGB}{70, 80, 95}    % Slate Blue-Grey
\definecolor{accent}{RGB}{220, 100, 50}     % Terracotta Accent
\definecolor{bodytext}{RGB}{40, 45, 55}     % Off-black for readability

% Basic page layout and paragraph styling
\setlength{\parskip}{0.6em}
\setlength{\parindent}{0pt}
\renewcommand{\baselinestretch}{1.15}

% Hyperlink configuration
\hypersetup{
    colorlinks=true,
    linkcolor=primary,
    filecolor=accent,
    urlcolor=primary,
    pdftitle={StitchFlow AI - Hackathon Phase 1 Proposal},
    pdfauthor={StitchFlow AI Team},
    pdfsubject={Multi-Agent Capacity Optimization and Intelligent Procurement},
}

% Style Section and Subsection headers
\titleformat{\section}
  {\color{primary}\normalfont\Large\bfseries}
  {\thesection}{1em}{}[{\color{accent}\titlerule[0.8pt]}]

\titleformat{\subsection}
  {\color{secondary}\normalfont\large\bfseries}
  {\thesubsection}{1em}{}

% Adjust list spacings
\setlist[itemize]{noitemsep, topsep=4pt, partopsep=0pt, parsep=2pt, leftmargin=1.8em}
\setlist[enumerate]{noitemsep, topsep=4pt, partopsep=0pt, parsep=2pt, leftmargin=1.8em}

% Header and Footer setup
\pagestyle{fancy}
\fancyhf{}
\lhead{\color{secondary}\small \textbf{StitchFlow AI} --- Concept Submission}
\rhead{\color{secondary}\small Phase 1: Idea \& Feasibility}
\lfoot{\color{secondary}\small \textbf{HACKATHON: AUTOMATE OR DIE}}
\rfoot{\color{secondary}\small Page \thepage}
\renewcommand{\headrulewidth}{0.5pt}
\renewcommand{\footrulewidth}{0.5pt}

\begin{document}

% Title Block
\begin{center}
    {\color{primary}\Huge \textbf{StitchFlow AI}}\\[0.4em]
    {\color{secondary}\large \textbf{Multi-Agent Capacity Optimization \& Intelligent Sourcing for Apparel Importers}}\\[0.8em]
    \textit{Theme: Production Planning \& Supply Chain | Subfield: Scheduling \& Capacity Optimization}\\
    \textit{Target Market: Tunisian Apparel Importers \& B2B Distributors}
    \vspace{1.5em}
\end{center}

\hrule
\vspace{1em}

\section{Problem Statement}
In developing economies, specifically in the Tunisian retail market, apparel importers and distributors act as the critical backbone of the supply chain. These businesses import high volumes of garments from global manufacturing hubs (primarily Turkey and China) and distribute them to local retail boutiques, shopping malls, and independent vendors. Managing this supply chain is a high-stakes balancing act where "capacity" is defined by capital availability, shipping logistics, and warehouse throughput.

Currently, this sector faces severe capacity mismatches driven by three main structural bottlenecks:
\begin{enumerate}
    \item \textbf{Long Import Lead Times and Slow Adjustments:} Sourcing apparel from Istanbul or Guangzhou requires a lead time of 6 to 12 weeks, including shipping routes to the Port of Rades and customs clearance processes. If importers misjudge local fashion tastes, they are stuck with the wrong inventory for months.
    \item \textbf{High Overstock Rates and Capital Lockup:} Operating on manual, spreadsheet-based forecast models and human "gut feel," importers face an average end-of-season overstock rate of \textbf{30\% to 40\%}. This slow-moving stock ties up valuable warehouse space (typically 1,500--3,000 $m^2$) and locks up vital working capital (averaging \textbf{150,000 to 300,000 TND} per season). Importers are forced to liquidate this stock at deep discounts (often 70\%+ off), destroying margins.
    \item \textbf{Understocking and Lost Sales (Stockouts):} Conversely, when a style goes viral on local social media (e.g., TikTok or Instagram boutique channels), importers cannot react in time. The average stockout rate on high-demand trending items is estimated at \textbf{25\%}, representing tens of thousands of dinars in lost top-line revenue and idle distribution capacity.
    \item \textbf{Manual and Fragmented Workflow:} Procurement managers spend over \textbf{15 hours per week} manually browsing supplier WhatsApp catalogs, checking Instagram trends, calling local boutique owners, and updating spreadsheets, leading to delayed decisions and high operational error rates.
\end{enumerate}

\section{Proposed Solution}
\textbf{StitchFlow AI} is an autonomous Multi-Agent System (MAS) designed to orchestrate capacity allocation and automate procurement decisions for apparel importers. It replaces intuition-driven purchasing with a real-time, closed-loop system that aligns import capacity with predicted local demand.

The high-level concept relies on deploying four specialized, collaborative AI agents:
\begin{itemize}
    \item \textbf{Internal Capacity Agent (Inward):} Continuously monitors internal databases (Odoo, SQLite, or Excel sheets), tracking warehouse bin occupancy (cubic meters), inventory aging, current SKU sales velocities, and available purchasing capital.
    \item \textbf{External Trend Scraper Agent (Outward):} Scrapes regional social media platforms (primarily Tunisian Instagram boutiques and TikTok fashion hashtags) and parses international supplier catalogs (e.g., Turkish wholesalers in Laleli/Merter) to identify rising styles.
    \item \textbf{Logistics \& Risk Agent:} Tracks external logistics signals including shipping transit delays to the Port of Rades, customs clearance times, and currency fluctuations (TND to EUR/USD).
    \item \textbf{Orchestrator Agent:} Consolidates the findings from the other agents, runs constraint-based optimization models, and outputs exact, risk-adjusted procurement and capacity allocation recommendations.
\end{itemize}

\textbf{Value Proposition:} StitchFlow AI ensures that every cubic meter of warehouse space and every dinar of importing capital is allocated to the highest-yielding apparel items, maximizing warehouse turnover, reducing overstock liquidations, and eliminating stockouts.

\section{Innovation}
Unlike standard forecasting software that relies on historical ERP data, StitchFlow AI introduced a decoupled, collaborative multi-agent architecture:
\begin{itemize}
    \item \textbf{Unstructured Social-to-Inventory Correlation:} It bridges the gap between unstructured local social media sentiment (Tunisian dialect and French/Arabic text/images from Instagram) and structured internal ERP databases.
    \item \textbf{Dynamic Capacity-Aware Optimization:} Standard inventory replenishment systems assume static lead times and infinite warehouse space. StitchFlow AI dynamically models lead times and warehouse capacity constraints. It uses a mathematical optimization engine to adjust order sizes based on current logistics delays and actual storage limits.
    \item \textbf{Autonomous Negotiation Prep \& Purchase Generation:} The system does not just show a chart; it generates ready-to-send purchase orders (POs) and drafts automated emails/WhatsApp messages to suppliers in Istanbul/Guangzhou, drastically reducing the purchasing loop.
\end{itemize}

\section{Quantified Impact}
Based on historical modeling of Tunisian mid-sized apparel importers, we project the following quantified business impacts upon implementing StitchFlow AI:
\begin{table}[h!]
\centering
\color{bodytext}
\begin{tabular}{lll}
\toprule
\textbf{Key Metric} & \textbf{Before StitchFlow AI} & \textbf{Projected After StitchFlow AI} \\
\midrule
End-of-Season Overstock Rate & 35\% & 15\% (\textbf{57\% reduction}) \\
Trend SKU Stockout Rate & 25\% & 5\% (\textbf{80\% reduction}) \\
Average Inventory Turn Rate & 3.2 turns/year & 4.8 turns/year (\textbf{50\% increase}) \\
Procurement Planning Time & 15 hours/week & 2 hours/week (\textbf{86\% time saved}) \\
Warehouse Capital Lockup & 200,000 TND & 80,000 TND (\textbf{120,000 TND freed}) \\
\bottomrule
\end{tabular}
\caption{Projected business impact of StitchFlow AI implementation.}
\end{table}

\section{Feasibility \& Hackathon Roadmap}
Building a functional prototype within a tight 36-hour hackathon window is highly feasible. Our implementation plan leverages lightweight, mature python libraries (CrewAI/LangChain, Streamlit, and PuLP) and a structured development roadmap:

\begin{itemize}
    \item \textbf{Hour 0 -- 12: Core Agent Development \& Data Mocking:} Build and configure the \textit{Internal Capacity Agent} (parsing Odoo/ERP mock databases) and the B2B catalog-scraping \textit{Trend Scraper Agent} (parsing Tunisian boutique Instagram trends and Turkish wholesale catalogs).
    \item \textbf{Hour 12 -- 24: Logistics Modeling \& Multi-Agent Orchestration:} Build the \textit{Logistics Agent} (simulating customs and delay factors at the Port of Rades). Integrate all agents into a collaborative system using CrewAI. Implement the mathematical optimization engine in Python (\texttt{PuLP}) to solve capacity-constrained procurement scenarios.
    \item \textbf{Hour 24 -- 36: UI Development \& System Validation:} Build the Streamlit dashboard displaying agent logs, warehouse occupancy levels, and automated purchase orders. Validate the end-to-end system under logistics and trend volatility shocks.
\end{itemize}

\section{Originality}
Supply chain optimization is traditionally reserved for massive enterprise players in the US/EU (e.g., Zara, Walmart) who have access to millions of dollars in IT infrastructure. 

StitchFlow AI is original because it applies advanced multi-agent technology to a traditionally low-tech, fragmented market: apparel importers in developing nations (starting with Tunisia). It democratizes predictive procurement by using LLM agents to convert raw social media data and unstructured wholesale chats into enterprise-grade capacity decisions. It turns the local constraints of Tunisian businesses (reliance on Turkish/Chinese imports, maritime shipping to Rades, high inflation, and limited credit) from operational weaknesses into optimized strategic advantages.

\section{Target Audience and Use Cases}
The primary target audience comprises \textbf{small-to-medium apparel importers and B2B distributors in Tunisia} (managing 10+ retail boutique clients or running their own physical store networks in Tunis, Sfax, and Sousse).

\textbf{Core Use Cases:}
\begin{enumerate}
    \item \textbf{Weekly Sourcing Run:} The Procurement Manager opens the dashboard on Monday morning. The Orchestrator Agent presents a list of 5 Turkish B2B styles that are trending on local Instagram and fits within the remaining warehouse capacity. The manager reviews and clicks "Approve" to generate purchase orders.
    \item \textbf{Warehouse Space Alerts:} When the system detects a slow-moving import shipment that will arrive simultaneously with a seasonal replenishment, the Capacity Agent alerts the manager and drafts a discount campaign to liquidate older inventory, freeing up rack capacity.
    \item \textbf{Customs/Logistics Shock Management:} If a shipping strike or customs delay occurs at the Port of Rades, the Logistics Agent adjusts estimated arrival times, prompting the Orchestrator to reroute capital to air-freight a smaller batch of high-margin items to prevent stockouts.
\end{enumerate}

\section{Competitive Differentiation}
StitchFlow AI represents a significant leap forward compared to existing industry alternatives:

\begin{enumerate}
    \item \textbf{Manual Spreadsheets \& Gut-Feel (Current Status Quo):}
    \begin{itemize}
        \item \textit{Limitations:} Highly reactive, prone to human error, blind to rapid social media trend shifts, and requires 15+ hours/week.
        \item \textit{Our Advantage:} 100\% automated, predictive, runs 24/7, and takes under 2 hours/week to manage.
    \end{itemize}
    \item \textbf{Traditional ERP Forecasting Modules (e.g., Odoo, SAP):}
    \begin{itemize}
        \item \textit{Limitations:} Rely purely on historical internal sales data. In fashion, history does not predict the next viral trend. They cannot parse external social media, local dialects, or unstructured supplier sheets.
        \item \textit{Our Advantage:} Correlates external social sentiment and supplier catalogs with internal warehouse capacity in real time.
    \end{itemize}
    \item \textbf{Traditional Marketing Trend Forecasting Services (e.g., WGSN):}
    \begin{itemize}
        \item \textit{Limitations:} Extremely expensive (thousands of dollars/year), tailored for massive Western brands, and completely detached from a Tunisian importer's actual warehouse capacity, logistics realities, or budget.
        \item \textit{Our Advantage:} Locally tailored, cheap, and directly integrated with the importer's physical and financial constraints.
    \end{itemize}
\end{enumerate}

\end{document}
